\chapter*{Resumen}

La sensibilidad y la alta resolución exigidas en la fotónica moderna hacen del diseño manual de sensores de Resonancia de Plasmones Superficiales (SPR) un proceso iterativo, computacionalmente costoso y tendiente a óptimos locales. Esta propuesta de investigación aborda el diseño y la optimización de estos sensores mediante un enfoque integral de Inteligencia Artificial ("AI for Science"), superando las limitaciones de los métodos tradicionales, como el Método de Matrices de Transferencia (TMM). El objetivo general consiste en desarrollar e implementar un entorno computacional inteligente que combine cuatro enfoques de vanguardia: (1) Diseño Inverso a través de redes neuronales profundas (DNN) para predecir automáticamente configuraciones óptimas de capas de materiales, (2) Modelos Sustitutos para acelerar el cálculo de respuestas electromagnéticas en tres órdenes de magnitud, (3) Redes Neuronales Informadas por la Física (PINNs) para garantizar el respeto de las condiciones de contorno de Maxwell minimizando la dependencia masiva de datos y (4) Inteligencia Artificial Explicable (XAI) para revelar correlaciones ocultas entre los espesores metálicos y el ancho de resonancia. Metodológicamente, se utilizará un marco en Python acelerado por GPU (NVIDIA RTX 4060 Ti) que alimentará un bucle optimizado de entrenamiento-verificación. El impacto esperado es una reducción dramática en el tiempo de diseño de biosensores ultradelgados, aportando una plataforma open-source validada y completamente reproducible que establezca un puente innovador entre el Aprendizaje Automático Avanzado y la Nanofotónica de precisión.


{\large \textbf{Palabras Clave} : }Inteligencia Artificial, Sensores Plasmónicos, Diseño Inverso, Aprendizaje Profundo, Modelos Sustitutos, PINNs, XAI.
